
\section{1. Introduction}

\subsection{1.1 Motivation}

Standard self-attention (Vaswani et al., 2017) allows every position to attend 
to every other position, requiring $\mathcal{O}(L^2)$ time and memory in 
sequence length. Empirically, however, attention is often \textbf{local}: for many 
heads and many tasks, most of the attention mass lies within a relatively 
small neighborhood around each query.

This observation motivates \textbf{window attention}: restrict each query to a 
window of size $w$ around it. Longformer (Beltagy et al., 2020), Mistral, and 
many modern efficient transformers use sliding windows with $w \in \{256, 
512, 1024\}$. Window attention reduces complexity to $\mathcal{O}(L w)$.

\textbf{The problem with fixed windows.} The window size is a static hyperparameter. 
But not all attention heads need the same window:

\begin{itemize}
\item Some heads are \textbf{local}: they attend to nearby tokens for syntax.
\item Some heads are \textbf{global}: they attend to distant tokens for coreference.
\item Some heads are \textbf{sparse}: they attend to a few specific positions.
\end{itemize}

A single window size forces all heads into the same trade-off. \textbf{Figure 1.1} 
shows the attention distance distribution for a middle GPT-2 layer: some heads 
concentrate within 32 tokens, others spread beyond 128.

\subsection{1.2 Problem Statement}

Given a transformer with $H$ heads and a maximum window size $S_\text{max}$, 
we seek a \textbf{per-head} window size $z_h \in [0, S_\text{max}]$ such that:

1. \textbf{Learned:} $z_h$ is trained end-to-end alongside the model.
2. \textbf{Differentiable:} gradients flow through $z_h$.
3. \textbf{Adaptive:} some heads converge to small $z_h$, others to large.
4. \textbf{Quality-preserving:} $\Delta\text{Acc} \leq 5\%$ on SST-2.

\subsection{1.3 Our Approach: Adaptive Span Attention}

We propose \textbf{Adaptive Span Attention} (ASA-Span), which learns per-head 
spans via:

1. \textbf{Continuous soft ramp.} Given span $z_h$ and ramp size $R$, the masking 
   function is:
\[
   g_{z_h}(d) = \text{clamp}\!\left(\frac{R + z_h - d}{R}, 0, 1\right)
   \tag{1.1}
\]
   where $d$ is the query-key distance. This is \textbf{piecewise linear} and 
   \textbf{differentiable} in $z_h$.

2. \textbf{Learnable span.} Each head maintains a scalar logit $\ell_h \in \mathbb{R}$:
\[
   z_h = \sigma(\ell_h) \cdot S_\text{max}
   \tag{1.2}
\]

3. \textbf{L1 span penalty.} A regularization term $\lambda \sum_h z_h$ encourages 
   the model to \textbf{shrink} spans that are not needed. At initialization, 
   $z_h \approx S_\text{max}$ (fully open); during training, spans shrink 
   to the smallest size preserving accuracy.

4. \textbf{No hard truncation.} Unlike sliding-window attention, ASA-Span does not 
   hard-truncate scores. It applies a soft multiplicative mask, so distant 
   tokens are \textbf{down-weighted} rather than removed. This preserves 
   differentiability and avoids artifacts at the window boundary.

\subsection{1.4 Contributions}

\begin{itemize}
\item \textbf{A differentiable soft-ramp formulation.} We extend Sukhbaatar et al.'s
\end{itemize}
  (2019) adaptive span to modern transformers with a numerically stable and 
  GPU-efficient implementation.
\begin{itemize}
\item \textbf{Per-head span heterogeneity.} We show that learned spans are highly
\end{itemize}
  non-uniform, with systematic dependence on head index and layer depth.
\begin{itemize}
\item \textbf{A production-hardened implementation.} Supports flat, nested, and causal
\end{itemize}
  architectures with RoPE, GQA, and T5-style relative position bias.
\begin{itemize}
\item \textbf{A systematic evaluation.} Benchmarked on DistilBERT, RoBERTa, and GPT-2
\end{itemize}
  on SST-2 with ablations over max span, ramp size, and L1 penalty.
\begin{itemize}
\item \textbf{Full reproducibility.} Code, checkpoints, and evaluation scripts released
\end{itemize}
  under MIT.

\subsection{1.5 Paper Organization}

Section 2 reviews related work. Section 3 provides background on attention and 
window attention. Section 4 presents ASA-Span. Section 5 describes integration 
with three architectures. Section 6 details the experimental setup. Section 
7 reports results. Section 8 analyzes the head heterogeneity and span 
distribution. Section 9 discusses practical guidance. Section 10 lists 
limitations. Section 11 concludes.



\subsection{Notation}
